\documentclass[11pt]{article}

\setlength{\topmargin}{0.0in}     % top of paper to head (less one inch)
\setlength{\headheight}{0in}    % height of the head
\setlength{\headsep}{0in}       % head to the top of the body
%\setlength{\textheight}{9.5in} % height of the body
\setlength{\textheight}{9in} % height of the body
\setlength{\oddsidemargin}{0mm} % left edge of paper to body (less one inch)
\setlength{\evensidemargin}{0mm} % ditto, even pages
\setlength{\textwidth}{6.5in}   % width of body
\setlength{\topskip}{0in}       % top of body to bottom of first line of text
%\setlength{\footheight}{0.25in} % height of footer
\setlength{\footskip}{0.50in}   % bottom of text to bottom of foot
% \renewcommand{\baselinestretch}{1.4}   % uncomment to double space



\newtheorem{theorem}{Theorem}
\newtheorem{corollary}{Corollary}
\newtheorem{lemma}{Lemma}
\newtheorem{observation}{Observation}
\newtheorem{definition}{Definition}
\newtheorem{fact}{Fact}
\newcommand{\proof}{\vspace*{-1ex} \noindent {\bf Proof: }}
\newcommand{\proofsketch}{\vspace*{-1ex} \noindent {\bf Proof Sketch: }}
\newcommand{\qed}{\hfill\rule{2mm}{2mm}}
\newcommand{\ceiling}[1]{{\left\lceil{#1}\right\rceil}}
\newcommand{\floor}[1]{{\left\lfloor{#1}\right\rfloor}}
\newcommand{\paren}[1]{\left({#1}\right)}
\newcommand{\braces}[1]{\left\{{#1}\right\}}
\newcommand{\brackets}[1]{\left[{#1}\right]}
\newcommand{\Prob}{{\rm Prob}}
\newcommand{\prob}{{\rm Prob}}
%
\newcommand{\ob}{$\langle$}
\newcommand{\cb}{$\rangle$}

\newcommand{\header}[3]{
   \pagestyle{plain}
   \noindent
   \begin{center}
   \framebox{
      \vbox{
    \hbox to 6.28in { {\bf CS 334 Computer Security
\hfill Fall 2014} }
       \vspace{4mm}
       \hbox to 6.28in { {\Large \hfill #1 \hfill} }
       \vspace{4mm}
    \hbox to 6.28in { {\sl #2 \hfill #3} }
      }
   }

   \end{center}
   \vspace*{4mm}
}

%\input{psfig.tex}

\begin{document}

\header{Handout: Introduction to Computer Security}{Assigned:
August 27}

\textit{Thanks to my UC-Berkeley colleagues for providing the basis for this 
handout.}

\section{The Scope of this Class}

My goal in this class is to teach you
the some of the most important and useful ideas in computer
security. By the end of this course, we hope you will have learned:
\begin{itemize}
\item {\emph How to build secure systems.}  You'll learn techniques for
designing, implementing, and maintaining secure systems.
\item {\emph How to evaluate the security of systems.}  Suppose someone hands you a
system they built. How do you tell whether their system is any good?
We'll discuss how systems have failed in the past, how attackers
break into systems in real life, and how to tell whether a given
system is likely to be secure.
\item {\emph How to communicate securely.} We'll cover topics from the science of
cryptography, which studies how several parties can communicate
securely over an insecure communications medium. 
\end{itemize}

Computer security is a broad field, that touches on almost every aspect of
computer science.  Moreover, given the prevalence of computers and computer networks
in industry, the subject is relevant to almost every aspect of commerce. 
I hope you'll enjoy the scenery as we explore!

What is computer security? Computer security is about computing
in the presence of an adversary. One might say that the defining
characteristic of the field, the lead character in the play, is the
adversary. Reliability, robustness, and fault tolerance are about how
to deal with Mother Nature, with random failures; in contrast,
security is about dealing with actions instigated by a knowledgeable
attacker who is dedicated to causing you harm. Security is about
surviving malice, and not just mischance. Wherever there is an
adversary, there is a computer security problem. 

Adversaries are all around us. The Code Red worm infected a quarter of
a million computers in less than a week, and contained a time-bomb set
to try to take down the White House web server on a specific
date. Fortunately, the attack on the White House was diverted, but one
research company is estimating the worm cost \$2 billion in lost
productivity and in cleaning up the mess caused by infected machines.
The Ponemon Institute, a cybersecurity research institute, estimated
that cyberinsecurity cost businesses over \$130 billion in
2011. Around Thanksgiving 2013, in the biggest retail hack in
U.S. history, Ukrainian hackers installed malware on Target stores
security and payment system that instructed the system to steal every
credit card used at all 1797 Target stores.  Since the malware was
running during the Christmas buying season, they managed to steal more
than 40 million credit card numbers.  This despite the fact that
Target was prepared for such an attack, having 6 months earlier
installed a 1.6 million dollar malware detection tool made by FireEye
(which DID detect the activity).  Currently many millions of computers
worldwide have been penetrated and ``owned'' by malicious parties;
many are used to send massive amounts of spam or make money through
phishing and identity fraud.  A 2010 Microsoft report suggested that
more than 11 million computers had been owned as of 2009 in the United
States alone.  The reality is that it's far more likely than not that
your laptop is already controlled by one of the major worldwide
botnets.  In effect, whatever information you store on it is already
available to, most likely, some large organized crime syndicate.  It's
a racket, and it pays well --- the perpetrators are raking in money
fast enough that they don't need a day job. How are we supposed to
secure our machines when there are folks like this out there? That's
the subject of this class.

\section{Risk Analysis}
If you remember nothing else from this class, remember that \textit{\textbf{computer security is essentially
risk analysis.}}  We will look at techniques that attempt to prevent data damage and theft from adversaries, 
but the reality is that absolute prevention is not possible, any more than it is possible to guarantee that
no one can break into your house.  Continuing with the house analogy, regardless of the number and types of locks and alarm systems you install, 
an adversary with sufficient technical skill and resources will be able to bypass any system you throw at her.
Even the extreme measure of hiring your
own private army to defend the house is no guarantee that unlawful entry will be prevented.  What each subsequent layer of
security does accomplish, however, is to raise the bar for the adversary.  Installing an alarm system will prevent some subset of
potential adversaries from entering your home.  Hiring an army will prevent many more (but not all) adversaries from entering your home.  
If you're measures are effective, each one increases the cost to the adversary (and likely increases your
costs as well).  Is the cost worthwhile?  That depends on the value you place on what you are trying to project.  Hiring an army
is not a cost effective means of preventing theft of your X-Box.  

Returning to computer security, the game, as it were, is to raise the bar high enough that the cost of the attack for the adversary
is higher than the value of the information they wish to compromise.  It's risk analysis, pure and simple:  the value of data, the
likelihood of perceived threats, the expertise and resources available to the adversary provide a measure of the risk.  Security
measures are then taken to lower that risk to an acceptable level.

If this reality is dissappointing, consider the problem from another
perspective.  The job of protecting a computer system is inherently
far more difficult than the job of attacking it, since system
defenders must thwart every possible attack (including those that may
never have been attempted before), while adversarys need only find a
single successful attack!  Given this inequity, it should not be
surprising that is it not realistic to successfully defend against all
possible attacks.

\section{It's all about the adversary}

The early history of computer security is interwoven with military applications
(probably because the military were one of the first big users of
computers, and the first to worry seriously about the potential for
misuse), so it should not be surprising that much of the terminology
has military connotations. We speak of an attacker who is trying to
attack computer systems, of defenders working to protect their system
from these threats, and so on. Well, you get the idea. 

It might be
surprising that we are going to spend so much time studying attackers
and thinking about how to break into systems. Aren't the attackers the
bad guys? Why on earth would we want to spread knowledge that will
help bad guys be more effective? 

Part of the answer is that you can't properly defend your system
unless you know how it might be attacked.  Civil engineers need to learn what makes bridges fall
down if they want to have any chance of building a bridge that will
remain standing. Software engineering is no different; you need to know
how systems fail in real life, if you want to have the best chance of
building a system that will resist attack. This means you'd better
know what kinds of attacks you are likely to face in the field. 

Understanding existing attacks is necessary, but not sufficient. 
Attackers are intelligent (well,
some of them are, but in this course we assume all of them are).  
If you deploy a new defense, they will
respond. If you build a new system, they will try to find its weak
points and attack there. Attackers adapt. This means that we have to
find ways to anticipate what kinds of attacks might be mounted against
us in the future. 

Security is like a game of chess, only it is one
where the attackers often get the last move. We design a system, and
then it is very hard to change once it has been deployed. If attackers
find a security hole in a widely deployed system, the consequences can
be pretty serious. Therefore, we'd better learn to predict in advance
what the attackers might do to us, so that we can eliminate most (and hopefully
all) the exploitable security holes before the system is deployed. We have to practice
thinking like an attacker, so that we will know in advance how secure
the system is. 

So what happens if you fail to anticipate new attacks?
The cellphone industry knows the answer. In the 1980's, they designed
and deployed an analog cellphone infrastructure with essentially no
security measures; cellphones transmitted all their billing
information in the clear, and security rested on the assumption that
attackers wouldn't bother to put together the equipment to intercept
it. That assumption held for a while, but sooner or later criminals
were bound to catch on, and they did. Technically savvy attackers
built ``black boxes'' that intercepted the radio communications and
cloned phones, and criminals used these to make fraudulent calls 
en masse and to mount call-selling operations for profit. Cellphone
operators were unprepared for this, and in the early 90's, it had
gotten so bad that the US cellphone carriers were losing more than \$1
billion per year. At one point almost 70\% of the
long-distance cellphone calls placed from downtown Oakland on a Friday
night were fraudulent. By this point the cellphone service providers
were already well aware that they had a serious problem, but because
it takes 5-10 years and a great deal of capital to replace the
deployed infrastructure of cellular base stations, they were in a
difficult position. This illustrates how failing to anticipate how your
system might be attacked, or underestimating the threat, can be a costly
mistake.

It is for these reasons that security design requires the
study of attacks. Security experts spend a lot of time trying to come
up with new attacks. This might sound counter-productive (why help the
attackers?), but it makes sense when you realize that it is better to
learn about vulnerabilities before the system is deployed than
after. If you know about the possible attacks in advance, you can
design a system to resist those attacks; anything else is a roll of
the dice. 

\section{A process for security evaluation}

How do we think
about the ways that an adversary might use to penetrate system
security or otherwise cause mischief? In this lecture, we're going to
develop a framework to help you think through these issues. 

The first
place to start, when analyzing a system, is its {\em security goals.} What
properties do we want the system to have, even when it is under
attack? What are we trying to protect from the attacker? Or, to look
at it the other way around, what are we trying to prevent?

Some common security goals: 
\begin{itemize}
\item {\em Confidentiality.} Often there is some
private information that we want to keep secret from the
adversary. Maybe it is a password, a bank account balance, or a diary
entry that we don't want anyone else to be able to read. It could be
anything. We want to prevent the adversary from learning our
secrets. 
\item {\em Integrity.} If the system stores some information, we might
  want to prevent the adversary from tampering with or modifying that
  information.  Integrity in the security context generally refers to
  the ``truth'' about a document or piece of data.  As an example,
  message integrity protocols seek to guarantee that the message that
  is received is the same as the message that was sent.  In the
  specific case that integrity refers to the sender or originator of
  some data, security researchers use the term {\em authentication}.
\item {\em Availability.} If the system performs some
function, it should be operational when we need it. Consequently, we
may need to prevent the adversary from taking the system out of
service at an inconvenient time.
\end{itemize}

For example, consider the
database of grade information that we use in this class. One obvious
goal is to protect its integrity, so that you can't just give yourself
an A+ merely by tampering with the grade database. University rules
require us to protect its confidentiality, so that no one else can
learn what grade you are getting. We probably also want some level of
availability, so that when the end of the semester comes we can
calculate the grades everyone will receive.  We might also be concerned with
authenticating the source of the data, so that we know that the grades were
created by someone who is authorized to do so.

Security goals can be
simple, or they can be detailed. Figuring out the set of security
goals that must be preserved is an exercise in requirements
analysis, they are the specification of what it means for a system to be
secure. The security goals are the goals we want to be met even when
an adversary is trying to violate them. You can recognize which goals
are security goals by asking yourself: if someone were to figure out
how to violate this goal, would it be considered a security breach? If
the answer is yes, you've found yourself a security goal. 

Security
goals are highly application-dependent, so it's hard to say much
more. Instead, I'll leave you with a famous quote from Young, Boebert,
and Kain: ``A program that has not been specified cannot be incorrect;
it can only be surprising.'' A system without security goals has not
been specified, and cannot be wrong; it can only be surprising. 

After
you have a set of security goals, the next step is to perform a 
{\em threat assessment}, which asks several questions. What kind of threats might
we face? What kind of capabilities might we expect adversaries to
have? What are the limits on what the adversary might be able to do to
us? The result is a threat model,a characterization of the threats the
system must deal with. 

When performing a threat assessment, we have
to decide how much we can predict about what kind of adversaries we
will be facing. Sometimes, we know very well who the adversary is, and
we may even know their capabilities, motivations, and limitations. For
instance, in the Cold War, the US military was oriented towards its
main enemy, the Soviets, and a lot of effort was put into
understanding the military capabilities of the USSR (how many
battalions of infantry do they have? how effective are their tanks?
how quickly can their navy respond to such-and-such threat?). When we
know what adversary we will be facing, we can craft a threat model
using that knowledge, so that our threat model reflects what that
particular adversary is likely to do to us and nothing
more.

However, all too often the adversary is not known. In this
case, we need to reason more generically about unavoidable limitations
that will be placed upon the adversary. As a light-hearted example,
physics tells us that the adversary can't go faster than the speed of
light, I don't care who they are, they can't violate the rules of
physics. That might be useful to know. More usefully, we can usually
look at the design of the system and identify what things an adversary
might be in a position to do. For instance, if the system is designed
so that secret information is never sent over a wireless network, then
we don't need to worry about the threat of eavesdropping upon the
wireless communications. If our system design is such that people
might discuss our secrets by phone, we had better include in our
threat model the possibility that an insider at the phone company
might be able to eavesdrop on our phone calls, or re-route them to the
wrong place, or fool people into thinking they are talking with
someone legitimate when actually they are speaking with the
attacker. 

A good threat model also specifies what threats we do not
care to defend against. For instance, if I want to analyze the
security of my home against burglary, I am not going to worry about
the threat that a team of burglars might fly a helicopter over my house
and rappel down my chimney to get into the house, Mission Impossible
style. There are far easier ways to break into my house, without going
to all that trouble. 

One can often classify adversaries according to
their motivation. For instance, consider adversaries who are motivated
by financial gain. It's a pretty safe bet that a financially-motivated
adversary is not going to spend more money on the attack than they
stand to gain from it. For instance, no burglar is likely to spend
thousands of dollars to steal my car radio; my car radio is simply not
worth that much. In general, motives are as varied as human nature,
and it is a good idea to be prepared for all eventualities. 

It's
often very helpful to look at the incentives of the various
parties. This is probably a familiar principle. Does the local fast
food joint make more profit on soft drinks than on the food? Then one
might expect some fast food places to take steps to boost sales of
soft drinks, perhaps salting its french fries heavily. Do customer
service representatives make a bonus if they handle more than a
certain number of calls per hour? Then one might expect some
representatives to be tempted to cut lengthy service calls short, or
to transfer trouble customers to other departments when possible. Do
spammers make money from everyone who responds to the spam, while
losing nothing from those who didn't wish to receive the spam? Then
one can expect that some spammers might be inclined to send their
emails as widely as possible, no matter how unpopular it makes
them. As a rule of thumb, organizations tend not to act against their
own self-interest, at least not too often. Incentives influence
behavior, not always, of course, but frequently enough to help
illuminate the motivations of potential adversaries. 

Incentives are
particularly relevant when two parties have opposing interests. When
incentives clash, conflict often follows. In this case it is worth
looking deeply at the potential for attacks by one such party against
the other. 

If threat assessment sounds difficult, just remember the
three W's: Who? How? Why? In other words: Who are the adversaries we
might face? How might they try to attack us, and what are their
capabilities? Why might they be motivated to attack us, and what are
their incentives? 

Finally, once we have the security goals and a
threat model, the last step is to perform a security analysis. The
goal of the security analysis is to see whether there is any attack
encompassed within the threat model that can successfully violate the
security goals. Security analysis is often highly technical and
depends on the details of the particular system being analyzed. We
will show you many methods for security analysis in the rest of the
course, without saying more now. 

An analogy: The security goals and
threat model defines the game, and the security analysis amounts to
figuring out who can win the game. The threat model defines the set of
moves the adversary is allowed to make, and the design of the system
defines how the defender will play the game. The security goals define
the success condition: if the adversary violates any security goal, he
wins; otherwise, the defender wins. The security analysis involves
examining all moves and counter-moves to see who has a winning
strategy. 

Another analogy: Mystery writers like to talk about means,
motive, and opportunity. That's not a bad way of thinking about what
we do during a security evaluation. The threat assessment examines the
means and motive; the security analysis examines what opportunity the
adversary might have to do harm. 

To sum up, evaluating the security
of a system involves three steps:
\begin{itemize}
\item {\em Identify the security goals.} What are we trying to protect?
\item {\em Perform a threat assessment.} What threats does the system need to protect against?
\item {\em Do a security analysis.} Can we envision any feasible attack that
would violate the security goals? This is the place where it can get
pretty technical.
\end{itemize}

The same process can be used for designing new
systems. The security goals and threat assessment is especially
relevant to system design, since it is usually easier to ensure
security when you know what security goals you are trying to achieve
and what threats you must protect against. And, as we perform our
security analysis, we can refine the system design to defend against
each new attack we discover.

\section{An example}
Enough slogans. Let's
do an example. Let's analyze the security of my home against
intruders.

What are my security goals? I'd like to protect the
assets in my home from theft or from being tampered with by
unauthorized parties (integrity). I'd like my family's safety to be
protected; for instance, if someone does break in to steal money, I'd
much prefer to know, so that no one surprises them and gets shot. I'd
like my house and its contents to remain in full working order
whenever I want them (availability).  I
sometimes want a certain measure of privacy when I'm home
(confidentiality). We could probably identify other security goals, but
that's more than enough to get us started.

Time for the threat
assessment. Who am I trying to protect against, and how might they be
motivated? A burglar might be motivated by financial gain. A peeping
Tom might be motivated by curiosity. Someone who holds a grudge
against me might want to secretly get revenge. And so on.

What capabilities (tools, skills, knowledge, access, etc.) might an
adversary have? What threats might I face? A burglar might have
lockpicks and the know-how to use them, or a crowbar to smash in the
door, or the capability to cut my phone lines. A repairman might have
unaccompanied access to the house for a time. Peepers might have
binoculars or a telescope. One of my neighbors might have
line-of-sight to my living room window, while another might be blocked
by trees. But there are probably some kinds of threats I'm willing to
ignore. For instance, I'm not worried that the fighter planes that 
occasionally fly over my neighborhood 
are going to start bombing my house. My house has no
effective way to defend against such an attack, but I doubt very much
that any officer dislikes me enough to throw away his career over
it. We could go on for pages, assessing which threats are realistic
given the possible motivations of the adversary, and ending up with a
detailed threat model. 

Finally, the security analysis. What kind of
attacks are possible? One can envision all sorts of crazy
scenarios. An everyday burglar might smash a window while I'm gone,
sneak in and grab some valuables, and run off before they're
caught. If I secure the windows, a determined burglar might take a
chainsaw to the walls and break in that way (believe it or not, it's
happened, though not to my house). 
A slightly smarter burglar might look under the flowerpot,
find the spare house key, and waltz right in. A really sneaky burglar
might throw a pebble against my window at 3am each morning, setting
off the burglar alarm and bringing the police running, for days at a
row, until the police decide to stop paying attention to my obviously
unreliable burglar alarm, and then the burglar can strike without
threat of police response. The neighbor with line-of-sight to my house
might use his telescope to peer in through my living window. Someone
intent on revenge might be able to leave unpleasant items on my lawn,
or, depending on my home's security system, might be able to smash my
windows without my noticing. A unscrupulous competitor who knows I
have an important business meeting with a potential client early
tomorrow morning might cut the power to my house in the middle of the
night, in hopes that my alarm clock will fail to ring, I will fail to
awaken, I will miss my meeting, and I will lose the client. One can
come up with some truly outlandish attack scenarios, but you probably
get the idea by now.

I hope this gives you some feeling for how a
security evaluation might go. If the example seems rather trivial, it
is probably because home security is already at least somewhat
familiar to you. However, when dealing with a complex computer system,
it helps to have a framework to structure the security evaluation, and
that's what the process outlined above is intended to help with.

\section{Terminology}

Lastly: Some terminology you might run into, and a
refresher on the terms we've been using so far. 
\begin{itemize}
\item {\em An attack} is an
attempt to breach system security. Not all attacks are
successful. 
\item {\em A threat} is a circumstance or scenario with the
potential to cause harm to a system. An attack usually refers to a
specific stratagem, whereas threat refers to a broader class of ways
that things could go wrong.
\item {\em A vulnerability} is an aspect of the
system that permits someone to mount a successful attack. Sometimes
called a {\em security hole}.  A {\em security weakness} is like a vulnerability,
only it is less clear whether it could actually lead to any direct
violation of the security goals. A weakness might represent a
potential vulnerability whose risk is unclear; or, several weaknesses
might combine to yield a full-fledged vulnerability.
\item {\em A security goal} is a goal that is supposed to be achieved by the system; if it
fails, the system will be considered insecure.
\item {\em A threat assessment} is an attempt to assess the set of all possible threats. A
threat model is a characterization of the possible threats, usually
produced during a threat assessment.
\item {\em 24/7} refers to the
window of time in which systems are most vulnerable to attack.  (OK,
this last one was a joke.)  
\end{itemize}

\end{document}
